\begin{abstract}
Four directions have been explored to strengthen multi-tool tool calling in LLM agents --- tool-aware fine-tuning, prompt scaffolding, KV-cache optimization, and activation steering --- each with a different trade-off between training cost and task-level targeting. This paper presents a training-free rotational steering method that belongs to the fourth category. The method consists of two history-free operators on a per-head ontology subspace $P_{\mathrm{ont}}$ derived directly from the tool catalog. The first is query-coverage $Q' = Q + \beta Q P_{\mathrm{ont}}$ ($\beta<0$), which contracts the query along the subspace; the second is a layer-adaptive K+Q that adds a K-side rotation on the first $\lfloor L/4\rfloor$ layers and keeps query-coverage across all layers. Unlike the stationary key-side amplifier shared by SEKA and AdaSEKA, both operators rotate attention \emph{away from the direction it already went}. Why this design matters for multi-tool is stated in one theorem: a stationary key-side operator cannot structurally encode which tools have been emitted (Theorem~\ref{thm:no-memory}), and the observable consequence is a rise in \texttt{repeated\_first\_tool\_rate}. Our two central claims are the following. On \emph{tool-selection performance} (\S\ref{sec:exp:performance}), under matched prompt/scorer/decoding protocol our operators yield consistent gains over SEKA and AdaSEKA on Qwen2.5-7B-Instruct for both MetaTool Subtask4 and $\tau^2$-bench retail; the gain is interpreted clearly not by aggregate F1 but by a position-level decomposition. On \emph{efficiency} (\S\ref{sec:exp:efficiency}), the method requires no weight training and no prompt expansion, working purely as inference-time low-rank hooks: training cost is zero relative to tool-aware fine-tuning, token overhead is zero relative to prompt scaffolding, and --- unlike KV-cache optimization methods --- it provides task-level selection improvement while staying compatible with them. Replacing the catalog-derived ontology basis with feature-shuffled, random, or PCA-of-K alternatives removes the effect, attributing the gain to the operator design rather than to basis quality.
\end{abstract}
